\documentclass[12pt,a4paper]{article}
\usepackage[spanish]{babel}
\usepackage[utf8]{inputenc}
\usepackage[T1]{fontenc}
\usepackage{lmodern}
\usepackage[a4paper,margin=2.4cm,headheight=14pt]{geometry}
\usepackage{xcolor}
\usepackage{fancyhdr}
\usepackage{amsmath,amssymb}
\usepackage{enumitem}
\usepackage{booktabs}
\usepackage{hyperref}
\usepackage{graphicx}

\definecolor{navy}{HTML}{0F2D8C}
\definecolor{accent}{HTML}{2563EB}
\definecolor{lightbg}{HTML}{EEF2FF}
\definecolor{answerbg}{HTML}{ECFDF5}
\definecolor{answerbd}{HTML}{10B981}

\hypersetup{colorlinks=true,linkcolor=navy,urlcolor=accent,citecolor=navy}

\pagestyle{fancy}
\fancyhf{}
\fancyhead[L]{\small\textcolor{navy}{\textbf{UNPRG — Lógica Simbólica}}}
\fancyhead[R]{\small\textcolor{gray}{Trabajo 01 · Lab 4}}
\fancyfoot[C]{\small\textcolor{gray}{\thepage}}
\renewcommand{\headrulewidth}{0.8pt}
\renewcommand{\headrule}{\hbox to\headwidth{\color{navy}\leaders\hrule height \headrulewidth\hfill}}

\newcommand{\respuesta}[1]{%
  \par\medskip\noindent
  \fcolorbox{answerbd}{answerbg}{%
    \parbox{\dimexpr\linewidth-2\fboxsep-2\fboxrule}{%
      \small\textbf{\textcolor{answerbd}{Respuesta:}} #1
    }%
  }\par\medskip
}

\newcommand{\datosbox}[1]{%
  \noindent\colorbox{lightbg}{%
    \parbox{\dimexpr\linewidth-2\fboxsep}{%
      \small #1
    }%
  }\par\medskip
}

\begin{document}

% ── PORTADA ────────────────────────────────────────────────────────
\begin{titlepage}
  \centering
  \vspace*{0.4cm}
  {\large\scshape\textcolor{navy}{Universidad Nacional Pedro Ruiz Gallo}\par}
  \vspace{0.15cm}
  {\small\textcolor{gray}{Facultad de Ingeniería Civil, Sistemas y Arquitectura}\par}
  \vspace{0.35cm}
  {\color{navy}\rule{\linewidth}{1.2pt}\par}
  \vspace{0.9cm}
  {\Huge\textbf{\textcolor{navy}{Lógica Simbólica}}\par}
  \vspace{0.35cm}
  {\Large\textcolor{accent}{Trabajo 01 — Resolución del Laboratorio 4}\par}
  \vspace{0.2cm}
  {\large\textcolor{gray}{Conjuntos y combinatoria}\par}
  \vspace{1.1cm}
  \begin{center}
    \begin{tabular}{c}
      \toprule
      \small\textcolor{navy}{\textbf{Integrantes}} \\ \midrule
      \small LUIS ENRIQUE CRISTOBAL DIAZ CAYACA \\
      \small MAURICIO MILLET FERNANDEZ LIZARZABURU \\
      \bottomrule
    \end{tabular}
  \end{center}
  \vspace{0.7cm}
  {\small\textbf{Ciclo:} II\par}
  \vspace{0.15cm}
  {\small\textbf{Profesor:} Dr. Mardo Victor Gonzales Herrera\par}
  \vspace{1.6cm}
  {\small\textcolor{gray}{26 de agosto de 2026 — Lambayeque, Perú}\par}
  \vfill
\end{titlepage}

% ── CONTENIDO ─────────────────────────────────────────────────────
\section*{\textcolor{navy}{Resolución de ejercicios}}
\subsection*{\textcolor{navy}{Conjuntos y combinatoria}}

% ── EJ 1 ──────────────────────────────────────────────────────────
\subsection*{\textcolor{accent}{1. Jugadores de fútbol, básquet y béisbol}}
\datosbox{%
  \textbf{Datos:} Fútbol: 48 jugadores. Básquet: 25 jugadores. Béisbol: 30 jugadores. Total: 68 jugadores. En los tres deportes: 6 jugadores.%
}
Sea $x$ la cantidad de jugadores que figuran exactamente en dos deportes y $y$ la cantidad que figura exactamente en un deporte.

Como hay 68 jugadores:
\[
  x + y + 6 = 68 \qquad\Longrightarrow\qquad x + y = 62 \tag{1}
\]

Contamos las participaciones:
\[
  48 + 25 + 30 = 103
\]
Los que están en un deporte cuentan una vez, los que están en dos cuentan dos veces y los que están en los tres cuentan tres veces:
\[
  y + 2x + 3(6) = 103 \;\Longrightarrow\; y + 2x + 18 = 103 \;\Longrightarrow\; y + 2x = 85 \tag{2}
\]

Restando (1) de (2):
\[
  x = 23
\]
\respuesta{23 jugadores figuran exactamente en dos deportes.}

% ── EJ 2 ──────────────────────────────────────────────────────────
\subsection*{\textcolor{accent}{2. Conjuntos $A$ y $B$}}
\datosbox{\textbf{Datos:} $|A \cup B| = 24,\; |A - B| = 10,\; |B - A| = 6$}
Usamos:
\[
  |A \cup B| = |A - B| + |B - A| + |A \cap B|
\]
Entonces:
\[
  24 = 10 + 6 + |A \cap B| \;\Longrightarrow\; |A \cap B| = 8
\]
Por tanto:
\[
  |A| = |A - B| + |A \cap B| = 10 + 8 = 18
\]
\[
  |B| = |B - A| + |A \cap B| = 6 + 8 = 14
\]
Calculamos:
\[
  5|A| - 4|B| = 5(18) - 4(14) = 90 - 56 = 34
\]
\respuesta{34.}

% ── EJ 3 ──────────────────────────────────────────────────────────
\subsection*{\textcolor{accent}{3. Conjuntos $A$ y $B$ definidos por comprensión}}
\[
  A = \{2x \mid x \in \mathbb{Z},\; 3 < x < 9\} \qquad
  B = \{3n \mid n \in \mathbb{Z},\; 5 < n < 10\}
\]
Los valores enteros de $x$ son $x = 4,5,6,7,8$. Por tanto:
\[
  A = \{8, 10, 12, 14, 16\}
\]
Los valores enteros de $n$ son $n = 6,7,8,9$. Por tanto:
\[
  B = \{18, 21, 24, 27\}
\]
No tienen elementos en común:
\[
  A \cap B = \varnothing \qquad |A \cap B| = 0
\]
Además $|A \cup B| = 5 + 4 = 9$. Entonces:
\[
  |A \cap B| + |A \cup B| = 0 + 9 = 9
\]
\respuesta{9.}

% ── EJ 4 ──────────────────────────────────────────────────────────
\subsection*{\textcolor{accent}{4. Escoger tres números del 1 al 9 sin consecutivos}}
Tenemos los números $1,2,3,4,5,6,7,8,9$. Queremos escoger 3 números sin que haya números consecutivos.

Usamos la fórmula:
\[
  \binom{n - r + 1}{r} \qquad\text{con } n=9,\; r=3
\]
\[
  \binom{9-3+1}{3} = \binom{7}{3} = \frac{7!}{3!\,4!} = \frac{7\cdot 6\cdot 5}{3\cdot 2\cdot 1} = 35
\]
\respuesta{35 formas.}

% ── EJ 5 ──────────────────────────────────────────────────────────
\subsection*{\textcolor{accent}{5. Palabras con las cinco vocales}}
Las vocales son $A, E, I, O, U$.

\subsubsection*{a) Palabras de longitud 4 sin repetir vocales}
Debemos escoger y ordenar 4 de las 5 vocales:
\[
  P(5,4) = \frac{5!}{(5-4)!} = 5! = 120
\]
\respuesta{120 palabras.}

\subsubsection*{b) Palabras de longitud 5 sin repetir vocales}
Se utilizan las cinco vocales:
\[
  P(5,5) = 5! = 5\cdot 4\cdot 3\cdot 2\cdot 1 = 120
\]
\respuesta{120 palabras.}

% ── EJ 6 ──────────────────────────────────────────────────────────
\subsection*{\textcolor{accent}{6. Comité de tres alumnos}}
En una clase hay 35 alumnos y se desea elegir un comité de 3 alumnos. Como el orden no importa, usamos combinación:
\[
  \binom{35}{3} = \frac{35!}{3!\,32!} = \frac{35\cdot 34\cdot 33}{3\cdot 2\cdot 1} = \frac{39270}{6} = 6545
\]
\respuesta{6545 comités diferentes.}

\end{document}
